% !TeX program = pdflatex
% ─────────────────────────────────────────────────────────────────────────────
%  papers/aranha.tex — A ARANHA ESTIGMÉRGICA: a multiplicidade geométrica.
%
%  Compila: pdflatex aranha.tex
%
%  AUTOCONTIDO por decisão do coordenador (20/08/2026): toda a matemática de
%  que precisa está DENTRO. Onde um irmão já dizia a mesma coisa, o enunciado
%  foi DUPLICADO e não referenciado — este documento lê-se sem abrir outro.
%  As únicas coisas de fora são o estilo da capa e os medidores citados.
% ─────────────────────────────────────────────────────────────────────────────
\documentclass[11pt,a4paper]{article}
\usepackage[utf8]{inputenc}\usepackage[T1]{fontenc}\usepackage[brazil]{babel}
\usepackage{amsmath,amssymb,amsthm}
\usepackage[margin=2.4cm]{geometry}
\usepackage{booktabs}\usepackage{xcolor}
\usepackage[hidelinks]{hyperref}

\definecolor{cinza}{gray}{0.35}
\newcommand{\code}[1]{\texttt{\small #1}}
\newcommand{\abs}[1]{\lvert #1\rvert}
\newcommand{\N}{\ensuremath{\mathbb{N}}}
\newcommand{\Z}{\ensuremath{\mathbb{Z}}}
\newcommand{\Q}{\ensuremath{\mathbb{Q}}}
\newcommand{\R}{\ensuremath{\mathbb{R}}}
\newcommand{\medido}[1]{\par\medskip\noindent{\small\color{cinza}\textbf{Medido.} #1}\par\medskip}
\newcommand{\sepcol}{\quad\penalty0$\big|$\quad\penalty0}
\newcommand{\colunas}[3]{%
\par\smallskip\noindent{\footnotesize\raggedright
\textbf{prova}: #1\sepcol\textbf{realiza}: #2\sepcol\textbf{verifica}: #3\par}\smallskip}
\providecommand{\interfacen}[1]{}
\newtheorem{teorema}{Teorema}
\newtheorem{lema}[teorema]{Lema}
\newtheorem{definicao}[teorema]{Definição}
\newtheorem{corolario}[teorema]{Corolário}
\renewcommand{\proofname}{Demonstração}

% ─────────────────────────────────────────────────────────────────────────────
%  papers/gkcapa.tex — a capa do Reino, mínima, para os papers em `article`.
%
%  O estilo.tex completo é do `report` (títulos de capítulo, skak, …). Aqui fica
%  só o que a capa precisa, para o paper continuar a compilar sozinho com
%  pdflatex. O tradutor próprio (tests/tex) carrega o estilo.tex da raiz / o
%  symlink papers/estilo.tex e avalia o mesmo `\gkcapa`.
% ─────────────────────────────────────────────────────────────────────────────
\definecolor{ouro}{HTML}{B8912F}
\definecolor{ouroclaro}{HTML}{F3E9CE}
\definecolor{tinta}{HTML}{1E1B16}
\definecolor{regua}{HTML}{6E6858}
\providecommand{\gknota}{\fontsize{7.62}{11.05}\selectfont}
\providecommand{\gktexto}{\fontsize{10.50}{15.22}\selectfont}
\providecommand{\gksub}{\fontsize{12.33}{17.87}\selectfont}
\providecommand{\gksec}{\fontsize{14.47}{20.98}\selectfont}
\providecommand{\gkcap}{\fontsize{16.99}{24.63}\selectfont}
\providecommand{\gktit}{\fontsize{23.42}{33.95}\selectfont}
\providecommand{\Prj}{\mathbb{P}}
\providecommand{\Trf}{\mathcal{T}}
\providecommand{\gkcapa}[3]{%
\title{\vspace{-0.6cm}
{\color{ouro}\rule{\textwidth}{1.4pt}}\\[6mm]
\textbf{\gktit\color{tinta} Reino Dourado}\\[3mm]{\gksec\itshape\color{regua} Golden Kingdom}\\[8mm]
{\gkcap\scshape\color{ouro} #1}\\[5mm]
{\gksub\color{regua} #2}\\[2mm]
{\gktexto\color{regua}\itshape #3}\\[7mm]
{\color{ouro}\rule{\textwidth}{0.6pt}}\\[9mm]{\gknota\color{regua}\begin{minipage}{\textwidth}\setlength{\parindent}{0pt}%
\textbf{Aviso de Isenção Algorítmica:} O universo, as leis e as entidades contidas nestas páginas ---
incluindo felinos matriciais, aves de rapina algébricas e amuletos de controle de fluxo --- são
construções de pura ficção formal. Esta obra não descreve a realidade física, não serve como
engenharia reversa do cosmos e não constitui doutrina religiosa. Qualquer semelhança com bugs reais do seu
sistema operacional é mera coincidência (ou falta de reza). Prossiga por sua própria conta e
risco.\end{minipage}}}
\author{}\date{}}


\begin{document}
\interfacen{6}
\gkcapa{A Aranha}
       {multiplicidade geométrica · memória no espaço}
       {a dobra que a projecção faz, e a coordenada que a desfaz}
\maketitle

\begin{abstract}\noindent
Uma trajetória discreta realizada num espaço identifica índices distintos na
mesma célula. Essa identificação \emph{é} a dobra da trajetória, e o número de
índices por célula --- a \textbf{multiplicidade geométrica} $G$ --- é a memória
de que essa dobra aconteceu. Prova-se aqui que $G$ é registável por
\textbf{escrita local no espaço}, sem que o agente conserve a sua própria
história; que a dimensão do espaço entra num sítio só, a vizinhança, e como
$\abs{V(x)}=2n$; que a fibra perdida pela projeção volta como \emph{uma}
coordenada de folha, seja qual for a dimensão, com volta exata; que em espaços
ultramétricos a célula \emph{é} a bola e as bolas particionam; e que numa órbita
determinista $G$ conta o período. Fecha-se com a curva de Heighway obtida por
dois caminhos disjuntos --- uma régua de bits e a dobra $D_{k+1}=D_k+D_k^{*}$
--- e com a demonstração de que coincidem.

\medskip\noindent\textbf{Auto-contenção.} Todo enunciado usado está enunciado e
demonstrado neste documento. Onde um irmão da mesma obra já dizia o mesmo, o
texto foi duplicado em vez de referenciado.
\end{abstract}

% ═══════════════════════════════════════════════════════════════════════════
\section{O objeto}\label{sec:objeto}

\begin{definicao}[trajetória, realização, multiplicidade]\label{def:objeto}
Seja $I=\{0,1,\dots,N\}$ um segmento inicial de \N\ --- a \textbf{parametrização}
--- e $X$ um conjunto com igualdade decidível. Uma \textbf{realização} é uma
aplicação
\[
\pi\colon I\longrightarrow X ,
\]
lida como ``o índice $i$ ocupa a célula $\pi(i)$''. A \textbf{multiplicidade
geométrica} de $\pi$ é
\[
G\colon X\to\N,\qquad G(x)=\bigl\lvert\pi^{-1}(x)\bigr\rvert
=\bigl\lvert\{i\in I:\pi(i)=x\}\bigr\rvert .
\]
O \textbf{suporte} de $G$ é $\operatorname{supp}G=\{x\in X: G(x)>0\}=\operatorname{im}\pi$.
\end{definicao}

\medskip\noindent\textbf{A letra.} $\pi$ é aqui a \textbf{projeção} --- a letra
padrão do fibrado, e a mesma do $\pi_k$ com que o Maestro desce um andar da
torre. O \emph{número} $\pi$ aparece na §\ref{sec:serie}, e aparece
\textbf{construído}: três somas alternadas independentes em aritmética inteira,
sem um dígito citado. São duas coisas com a mesma letra e a convenção fica dita
aqui, no primeiro uso; onde houver risco, escreve-se ``o número $\pi$''.

\medskip\noindent Note-se o que \emph{não} está na definição: nem estrutura algébrica em
$X$, nem métrica, nem dimensão. Só igualdade. Esta economia não é elegância --- é
o conteúdo do Teorema~\ref{thm:multiplicidade}, e é o que se paga no
Teorema~\ref{thm:dimensao}, onde a dimensão finalmente aparece e se vê que
aparece uma vez só.

\begin{definicao}[dobra]\label{def:dobra}
A relação $i\sim j \iff \pi(i)=\pi(j)$ é uma equivalência em $I$. Chama-se
\textbf{dobra} de $\pi$, e as suas classes são as fibras $\pi^{-1}(x)$.
\end{definicao}

% ═══════════════════════════════════════════════════════════════════════════
\section{A multiplicidade: quatro cláusulas, nenhuma sobre o espaço}

\begin{teorema}[multiplicidade geométrica]\label{thm:multiplicidade}
Com a Def.~\ref{def:objeto}, sejam $G_t(x)=\lvert\{i\le t:\pi(i)=x\}\rvert$ os
campos parciais. Então:
\begin{enumerate}\itemsep3pt
\item[(1)] \textbf{Duplicidade é dobra.} Para $i\neq j$ tem-se $\pi(i)=\pi(j)$
se e só se $i$ e $j$ pertencem à mesma classe de $\sim$. A dobra da
Def.~\ref{def:dobra} é a fibra de $\pi$ --- não uma metáfora dela.
\item[(2)] \textbf{$G$ regista a dobra.} $G(x)>1$ se e só se existem $i\neq j$
com $\pi(i)=\pi(j)=x$. Em particular $\pi$ é injetiva se e só se $G\le1$.
\item[(3)] \textbf{A memória não pertence ao agente.} Vale a recorrência local
\[
G_{t+1}(x)=G_t(x)+\mathbf{1}_{\{x=\pi(t+1)\}},\qquad G_{-1}\equiv0,
\]
e $G=G_N$. Cada passo altera $G$ num único ponto, e esse ponto é a célula
corrente. Nenhum passo consulta $\pi(0),\dots,\pi(t)$.
\item[(4)] \textbf{Conservação.} $\displaystyle\sum_{x\in X}G(x)=\abs{I}=N+1$, e
a soma é finita porque $\operatorname{supp}G$ tem no máximo $\abs{I}$ elementos.
\end{enumerate}
\end{teorema}

\begin{proof}
(1) e (2) são a definição de fibra e a sua cardinalidade.

(3) Fixe-se $x$. Por definição $G_{t+1}(x)=\lvert\{i\le t+1:\pi(i)=x\}\rvert$, e
o conjunto contado difere do de $G_t(x)$ exatamente pelo elemento $t+1$, que lá
está se e só se $\pi(t+1)=x$. Donde a recorrência. A cláusula ``num único ponto''
é a leitura da mesma igualdade: para $x\neq\pi(t+1)$ a indicatriz anula-se e
$G_{t+1}(x)=G_t(x)$. O termo direito só depende de $t+1$ através de $\pi(t+1)$,
que é o dado corrente: a história não é consultada.

(4) As fibras $\pi^{-1}(x)$, $x\in\operatorname{im}\pi$, particionam $I$, e a
cardinalidade é aditiva sobre partições finitas:
$\sum_x\lvert\pi^{-1}(x)\rvert=\lvert I\rvert$. Fora do suporte os termos são
nulos.
\end{proof}

\medskip\noindent\textbf{O que a prova usou.} Só a igualdade em $X$, para
decidir $x=\pi(t+1)$ e para formar as fibras. Nenhuma das quatro cláusulas
menciona a estrutura de $X$. É este o resultado: a memória estigmérgica não é um
fenómeno do plano.

\medskip\noindent A composição que o teorema exibe é
\[
\boxed{\ \text{trajetória}\ \xrightarrow{\ \pi\ }\ \text{célula}\
\xrightarrow{\ +1\ }\ \text{campo }G\ \xrightarrow{\ \text{leitura local}\ }\
\text{decisão}\ }
\]
e o agente é o que fecha o ciclo lendo o que ele próprio escreveu, sem guardar
que o escreveu.

% ═══════════════════════════════════════════════════════════════════════════
\section{O algoritmo}\label{sec:algoritmo}

\noindent A cláusula (3) do Teorema~\ref{thm:multiplicidade} \emph{é} um
algoritmo. Escrito:

\medskip\noindent\textbf{Algoritmo A (a aranha).}
\begin{enumerate}\itemsep2pt
\item Inicializar $G\equiv0$ (o campo vazio).
\item Para $t=0,\dots,N$: ler $x\leftarrow\pi(t)$; escrever $G(x)\leftarrow G(x)+1$.
\item \emph{(opcional, o agente)} antes de escrever, ler a vizinhança
$\mathcal{P}_t(x)=\bigl(G_t(y)\bigr)_{y\in V(x)}$ e decidir o passo seguinte
com base nela.
\end{enumerate}

\medskip\noindent\textbf{Estado.} O algoritmo mantém $G$ e o índice corrente.
Não mantém a sequência. Este é o ponto e não uma otimização: um agente que
guardasse $\pi(0),\dots,\pi(t)$ teria a informação toda e o campo seria
supérfluo.

\medskip\noindent\textbf{Representação de $G$.} O campo é uma aplicação
$\operatorname{supp}G\to\N$ com $\lvert\operatorname{supp}G\rvert\le\abs{I}$.
Armazená-lo como \emph{tabela sobre $X$} --- por exemplo um arranjo
$G[m]\times[m]$ quando $X=\Z^2$ --- é uma escolha de representação, não o
algoritmo: custa $\lvert X\rvert$ e não $\abs{I}$, e em $\Z^3$ com $m=2048$ pede
$2048^3$ células. A representação fiel à cláusula (3) é \textbf{esparsa}: a
memória existe onde a trajetória passou. Com endereçamento aberto num arranjo
fixo de capacidade $C\ge2\abs{I}$, o custo é $O(1)$ amortizado por passo e a
memória total é $O(\abs{I})$, decidida antes de correr --- não há alocação
dinâmica em parte nenhuma do algoritmo.

\medskip\noindent\textbf{Complexidade.} $N+1$ passos, cada um com uma leitura e
uma escrita: $\Theta(\abs{I})$ operações, $\Theta(\abs{I})$ memória, e nenhuma
dependência de $\lvert X\rvert$.

% ═══════════════════════════════════════════════════════════════════════════
\section{Onde a dimensão entra}\label{sec:dimensao}

\noindent A cláusula (3) do algoritmo pede uma \emph{vizinhança}. É aqui, e só
aqui, que a estrutura de $X$ é chamada.

\begin{definicao}[a grade e a sua vizinhança]\label{def:grade}
Em $X=\Z^n$ com a base canónica $e_1,\dots,e_n$, a vizinhança de $x$ é
\[
V(x)=\{x\pm e_i : i=1,\dots,n\}.
\]
\end{definicao}

\begin{teorema}[o $n$ entra num sítio só]\label{thm:dimensao}
Em $\Z^n$ tem-se $\lvert V(x)\rvert=2n$ para todo $x$, e os $2n$ vizinhos são
distintos. Além disso, o Teorema~\ref{thm:multiplicidade} e o Algoritmo~A valem
palavra por palavra em $\Z^n$ para todo $n\ge1$, sem alteração nenhuma: a
dimensão não ocorre em nenhuma das quatro cláusulas.
\end{teorema}

\begin{proof}
Os $2n$ pontos $x\pm e_i$ são distintos: $x+\varepsilon e_i=x+\varepsilon'e_j$
com $\varepsilon,\varepsilon'\in\{+1,-1\}$ dá $\varepsilon e_i=\varepsilon'e_j$,
e como os $e_i$ são livres isso força $i=j$ e $\varepsilon=\varepsilon'$. A
segunda afirmação é imediata da demonstração do
Teorema~\ref{thm:multiplicidade}, que usou apenas a igualdade em $X$, e $\Z^n$
tem igualdade decidível coordenada a coordenada.
\end{proof}

\medskip\noindent\textbf{Leitura.} Enunciar o teorema em $\Z^2$ e realizá-lo num
arranjo do plano dá a impressão de que o $2$ é do teorema. Não é: é do exemplo, e
do arranjo. O $n$ aparece uma vez, no tamanho da vizinhança.

% ═══════════════════════════════════════════════════════════════════════════
\section{O levantamento, e a volta}\label{sec:inversa}

\noindent A projeção $\pi$ perde exatamente a dobra. A pergunta dual: pode a
dobra ser \emph{desfeita} sem inventar geometria? Sim, e ao preço de uma
coordenada.

\begin{definicao}[número de visita]\label{def:k}
Para $i\in I$,
\[
k(i)=\bigl\lvert\{j\in I: j\le i,\ \pi(j)=\pi(i)\}\bigr\rvert .
\]
\end{definicao}

\begin{teorema}[levantamento em folhas]\label{thm:levantamento}
Seja $\tilde\pi\colon I\to X\times\N$, $\tilde\pi(i)=\bigl(\pi(i),k(i)\bigr)$, e
$\tilde G$ a multiplicidade de $\tilde\pi$. Então:
\begin{enumerate}\itemsep3pt
\item[(1)] \textbf{Injetividade:} $\tilde\pi$ é injetiva, e portanto
$\tilde G\equiv1$ em $\operatorname{im}\tilde\pi$.
\item[(2)] \textbf{Projeção:} $\operatorname{pr}_1\circ\tilde\pi=\pi$.
\item[(3)] \textbf{Fibra vertical:} para cada $x\in\operatorname{im}\pi$,
$\{k(i):\pi(i)=x\}=\{1,2,\dots,G(x)\}$, exatamente.
\item[(4)] \textbf{Conservação:}
$\lvert\operatorname{im}\tilde\pi\rvert=\abs{I}=\sum_x G(x)$.
\item[(5)] \textbf{Volta:} $G$ recupera-se de $\tilde\pi$ sem qualquer outro
dado, por $G(x)=\bigl\lvert\{i:\operatorname{pr}_1\tilde\pi(i)=x\}\bigr\rvert$.
\end{enumerate}
Além disso, se $X=\Z^n$ então $X\times\N\subseteq\Z^{n+1}$: o levantamento sobe
\textbf{um} andar, seja qual for $n$.
\end{teorema}

\begin{proof}
(3) primeiro, porque as outras se apoiam nela. Fixe-se $x$ e sejam
$i_1<i_2<\dots<i_g$ os índices da fibra $\pi^{-1}(x)$, $g=G(x)$. Por
definição $k(i_r)=\lvert\{j\le i_r:\pi(j)=x\}\rvert=\lvert\{i_1,\dots,i_r\}\rvert=r$.
Logo $k$ restrita à fibra é a enumeração $i_r\mapsto r$, que é uma bijeção sobre
$\{1,\dots,g\}$.

(1) Sejam $\tilde\pi(i)=\tilde\pi(j)$. Então $\pi(i)=\pi(j)=:x$, logo $i$ e $j$
estão na mesma fibra, e $k(i)=k(j)$. Por (3) $k$ é injetiva nessa fibra, donde
$i=j$. Que $\tilde G\equiv1$ é a cláusula (2) do
Teorema~\ref{thm:multiplicidade} aplicada a $\tilde\pi$.

(2) $\operatorname{pr}_1\bigl(\pi(i),k(i)\bigr)=\pi(i)$, por definição de
$\operatorname{pr}_1$.

(4) Por (1), $\tilde\pi$ é uma bijeção de $I$ sobre a sua imagem, donde
$\lvert\operatorname{im}\tilde\pi\rvert=\abs{I}$; a segunda igualdade é a
cláusula (4) do Teorema~\ref{thm:multiplicidade}.

(5) Por (2), $\{i:\operatorname{pr}_1\tilde\pi(i)=x\}=\{i:\pi(i)=x\}=\pi^{-1}(x)$,
cuja cardinalidade é $G(x)$ por definição.

A última afirmação é a inclusão $\Z^n\times\N\subseteq\Z^n\times\Z=\Z^{n+1}$.
\end{proof}

\medskip\noindent\textbf{Algoritmo B (a aranha inversa).}
\begin{enumerate}\itemsep2pt
\item Inicializar um campo $c\equiv0$ sobre $X$.
\item Para $i=0,\dots,N$: ler $x\leftarrow\pi(i)$; escrever
$c(x)\leftarrow c(x)+1$; emitir $\tilde\pi(i)=\bigl(x,\,c(x)\bigr)$.
\end{enumerate}
O campo $c$ é o próprio $G$ parcial: o Algoritmo~B é o Algoritmo~A com o valor
lido a sair. Não inventa geometria --- desfaz a identificação $i\sim j$ que a
realização tinha feito.

\medskip\noindent\textbf{Três observações que a prática obriga a escrever.}

\emph{Primeira:} a injetividade de $\tilde\pi$ não precisa de ser verificada par
a par. Basta correr o Algoritmo~A \emph{sobre $\tilde\pi$}, em dimensão $n+1$, e
ler $\max\tilde G$. Isto é uma verificação por outro caminho, não a mesma conta
duas vezes: a prova de (1) usa (3) e a ordem dos índices; a corrida usa só a
igualdade de vetores.

\emph{Segunda:} $\operatorname{pr}_1$ não é um passo que se execute. Correr o
campo em dimensão $n$ sobre os vetores de $\Z^{n+1}$ \emph{é} a projeção --- a
folha está no vetor e o campo não a lê. Escrever ``zerar a última coordenada''
antes de recontar é uma linha que não faz nada.

\emph{Terceira:} a volta (5) é o que distingue um \emph{inverso} de um
\emph{levantamento}. Um levantamento que não devolvesse $G$ teria acrescentado
informação em vez de reorganizar a que havia.

% ═══════════════════════════════════════════════════════════════════════════
\section{Espaços ultramétricos: a célula é a bola}\label{sec:ultra}

\noindent O Teorema~\ref{thm:multiplicidade} não pediu métrica. Mas quando $X$
tem uma, a relação entre $G$ e as bolas diz o que a métrica é.

\begin{definicao}[a árvore do encaixe]\label{def:arvore}
Para $a,b$ inteiros de $p$ bits, seja $\delta(a,b)$ a profundidade da primeira
divergência: o menor $q\ge0$ tal que os bits de peso $p-1,\dots,p-q$ de $a$ e $b$
coincidem e o de peso $p-1-q$ difere; e $\delta(a,a)=p$. Ponha-se
\[
d(a,b)=2^{-\delta(a,b)} .
\]
\end{definicao}

\begin{lema}[$d$ é ultramétrica]\label{lem:ultra}
$d(a,c)\le\max\{d(a,b),d(b,c)\}$ para quaisquer $a,b,c$.
\end{lema}

\begin{proof}
Se $a$ e $b$ concordam nos $\delta(a,b)$ bits mais altos e $b$ e $c$ nos
$\delta(b,c)$, então $a$ e $c$ concordam nos primeiros
$\min\{\delta(a,b),\delta(b,c)\}$, logo
$\delta(a,c)\ge\min\{\delta(a,b),\delta(b,c)\}$. Aplicando $t\mapsto2^{-t}$, que
inverte a ordem, vem $d(a,c)\le\max\{d(a,b),d(b,c)\}$.
\end{proof}

\begin{teorema}[as bolas particionam]\label{thm:bolas}
Seja $(X,d)$ ultramétrico e $B(a,r)=\{x: d(a,x)\le r\}$. Então:
\begin{enumerate}\itemsep3pt
\item[(1)] \textbf{Todo ponto é centro:} se $b\in B(a,r)$ então $B(b,r)=B(a,r)$.
\item[(2)] \textbf{Partição:} duas bolas de raio $r$ ou coincidem ou são
disjuntas. Em particular, para cada $r$ as bolas de raio $r$ formam uma partição
de $X$.
\item[(3)] Fixado $r$, seja $\pi_r(i)$ a bola de raio $r$ que contém $i$. Então
$\pi_r$ é uma realização no sentido da Def.~\ref{def:objeto}, a sua célula
\emph{é} a bola, e $G(x)$ é o número de índices dentro dela.
\end{enumerate}
\end{teorema}

\begin{proof}
(1) Seja $x\in B(a,r)$. Então $d(b,x)\le\max\{d(b,a),d(a,x)\}\le r$, logo
$x\in B(b,r)$ e $B(a,r)\subseteq B(b,r)$. Como $b\in B(a,r)$ implica
$a\in B(b,r)$ (a ultramétrica é simétrica), a inclusão contrária sai do mesmo
argumento com os papéis trocados.

(2) Se $c\in B(a,r)\cap B(b,r)$, por (1) tem-se $B(a,r)=B(c,r)=B(b,r)$. Se a
interseção é vazia, nada há a provar. Que cobrem $X$ é imediato:
$a\in B(a,r)$.

(3) A aplicação $i\mapsto\pi_r(i)$ está bem definida por (2) --- há
\emph{exatamente} uma bola de raio $r$ contendo $i$ --- e o contradomínio tem
igualdade decidível (duas bolas são iguais ou disjuntas, e testa-se pela
interseção). A afirmação sobre $G$ é a Def.~\ref{def:objeto}.
\end{proof}

\medskip\noindent\textbf{O que a ultramétrica é, e o que não é.} Não é uma
desigualdade triangular ``mais apertada'' de que se tire um número melhor: é a
propriedade de partição do Teorema~\ref{thm:bolas}. Numa métrica euclidiana as
bolas de raio fixo cruzam-se parcialmente --- existem $a,b$ e pontos $x,y,z$ com
$x\in B(a)\cap B(b)$, $y\in B(a)\setminus B(b)$, $z\in B(b)\setminus B(a)$ --- e
nenhum ponto interior é centro de todas as bolas que o contêm. É por isso que
$\sum_x G(x)=\abs{I}$ não exige cuidado nenhum aqui: não há sobreposição a haver.

\medskip\noindent\textbf{E é a métrica do corte.} Dois reais construídos como
cortes estão próximos nesta métrica quando \emph{concordam até fundo}: a
distância decai com a profundidade em que a expansão binária ainda coincide.
É a topologia que o encaixe induz, e o campo $G$ corre nela com o Algoritmo~A
inalterado --- a única coisa que $X$ tinha de ter era igualdade, e bolas
comparam-se por igualdade.

% ═══════════════════════════════════════════════════════════════════════════
\section{Órbitas deterministas: $G$ conta o período}\label{sec:periodo}

\begin{definicao}[realização determinista]\label{def:determinista}
$\pi$ diz-se \textbf{determinista} se existe $f\colon X\to X$ com
$\pi(t+1)=f(\pi(t))$ para todo $t<N$.
\end{definicao}

\begin{teorema}[o período lê-se no campo]\label{thm:periodo}
Seja $\pi$ determinista. Se existem $i<j$ com $\pi(i)=\pi(j)$, sejam
\[
\mu=\min\{i:\exists j>i,\ \pi(j)=\pi(i)\}\quad(\text{a \textbf{cauda}}),\qquad
p=\min\{j-\mu:\ j>\mu,\ \pi(j)=\pi(\mu)\}\quad(\text{o \textbf{período}}).
\]
Se $N\ge\mu+2p-1$, então
\[
\bigl\lvert\{x: G(x)>1\}\bigr\rvert=p
\qquad\text{e}\qquad
\bigl\lvert\{x: G(x)=1\}\bigr\rvert=\mu .
\]
Se, pelo contrário, $\pi$ é injetiva, então $G\equiv1$ e ambos os conjuntos são,
respetivamente, vazio e $\operatorname{im}\pi$.
\end{teorema}

\begin{proof}
Comece-se pela propriedade que o determinismo dá: se $\pi(i)=\pi(j)$ então
$\pi(i+m)=\pi(j+m)$ para todo $m$ com $j+m\le N$, por indução em $m$ aplicando
$f$. Logo a sucessão é periódica de período $p$ a partir de $\mu$:
$\pi(t+p)=\pi(t)$ para todo $t\ge\mu$ com $t+p\le N$.

\emph{As células do ciclo são $p$ e são distintas.} Sejam
$C=\{\pi(\mu),\pi(\mu+1),\dots,\pi(\mu+p-1)\}$. Se fosse
$\pi(\mu+r)=\pi(\mu+s)$ com $0\le r<s<p$, aplicando $f$ $\,(p-s)$ vezes viria
$\pi(\mu+r+p-s)=\pi(\mu+p)=\pi(\mu)$ com $0<r+p-s<p$, contra a minimalidade de
$p$. Logo $\lvert C\rvert=p$.

\emph{As células da cauda são $\mu$, distintas entre si e disjuntas de $C$.}
Sejam $T=\{\pi(0),\dots,\pi(\mu-1)\}$. Se $\pi(a)=\pi(b)$ com $a<b<\mu$, então
$a$ seria um índice repetido menor que $\mu$, contra a minimalidade de $\mu$;
logo $\lvert T\rvert=\mu$. Se $\pi(a)=\pi(\mu+r)$ com $a<\mu$ e $0\le r<p$,
aplicando $f$ $\,(p-r)$ vezes viria $\pi(a+p-r)=\pi(\mu)$ --- e como
$a<\mu$, o índice $a$ repete-se, contra a minimalidade de $\mu$ novamente.

\emph{Contagem.} Como $N\ge\mu+2p-1$, a órbita completa pelo menos duas voltas:
cada $x\in C$ é atingido em dois índices distintos $\mu+r$ e $\mu+r+p$, ambos
$\le N$, donde $G(x)\ge2$. Cada $x\in T$ é atingido uma só vez, por
$T\cap C=\emptyset$ e pela injetividade em $T$, donde $G(x)=1$. E
$\operatorname{im}\pi=T\cup C$. Logo $\{G>1\}=C$ com $p$ elementos e $\{G=1\}=T$
com $\mu$ elementos.

O caso injetivo é a cláusula (2) do Teorema~\ref{thm:multiplicidade}.
\end{proof}

\medskip\noindent\textbf{A hipótese $N\ge\mu+2p-1$ é necessária.} Sem ela o
enunciado é falso, e falso do modo mais banal: se a trajetória parar depois de
entrar no ciclo mas antes de o fechar duas vezes, as células do ciclo ainda
visitadas uma só vez contam para $\{G=1\}$. Exemplo mínimo, um passo abaixo do
limiar: $f(z)=z^2-1$ sobre \Z\ com $\pi(0)=0$ e $N=2$ dá a órbita $0,-1,0$, com
$\mu=0$ e $p=2$, logo o limiar é $N\ge3$; e
$\lvert\{G>1\}\rvert=1\neq2=p$. Com $N=3$ a órbita é $0,-1,0,-1$ e o enunciado
vale.

\medskip\noindent O contra-exemplo \emph{não} pode ser $N=1$. Aí a órbita é
$0,-1$ sem repetição nenhuma, a hipótese do próprio teorema --- existirem $i<j$
com $\pi(i)=\pi(j)$ --- falha, e $\mu$ e $p$ nem estão definidos pela
trajetória. Um caso onde a hipótese principal não se verifica não testemunha
nada sobre a hipótese auxiliar.

\medskip\noindent\textbf{Leitura: a deteção de ciclo sem o segundo ponteiro.} O
algoritmo clássico de deteção de ciclo mantém dois cursores a velocidades
diferentes precisamente porque o agente não tem memória do percurso. O
Teorema~\ref{thm:periodo} diz que, se a memória estiver \emph{no espaço}, um
cursor basta: $\mu$ e $p$ lêem-se no campo, sem segunda passagem e sem comparar
posições. É a cláusula (3) do Teorema~\ref{thm:multiplicidade} a pagar-se.

% ═══════════════════════════════════════════════════════════════════════════
\section{O que o passo lê: o espaço ou o tempo}\label{sec:leitura}

\noindent A Def.~\ref{def:determinista} pede que o passo seja função da célula.
Isto não separa objetos matemáticos de outros --- separa \textbf{duas
dependências}, e ambas são construções correntes. O teorema seguinte diz que a
escolha não é livre.

\begin{teorema}[dobrar sem período obriga a ler o índice]\label{thm:leitura}
Seja $\pi\colon I\to X$ com $G(x)>1$ para algum $x$ (dobra) e tal que $\pi$
\emph{não} é eventualmente periódica. Então não existe $f\colon X\to X$ com
$\pi(t+1)=f(\pi(t))$: o passo depende necessariamente do índice.
\end{teorema}

\begin{proof}
Suponha-se que $f$ existe. Como $\pi$ dobra, há $i<j$ com $\pi(i)=\pi(j)$.
Aplicando $f$ repetidamente, $\pi(i+m)=\pi(j+m)$ para todo $m$ admissível, logo
$\pi(t+(j-i))=\pi(t)$ para todo $t\ge i$: $\pi$ é eventualmente periódica de
período $j-i$ a partir de $i$, contra a hipótese.
\end{proof}

\begin{corolario}[a aranha e o dragão são duais]\label{cor:duais}
Sejam as duas maneiras de uma trajetória guardar a sua própria história:
\begin{center}\small
\begin{tabular}{@{}lll@{}}
\toprule
 & \textbf{a aranha} & \textbf{o dragão} \\
\midrule
a memória está & no \textbf{espaço} ($G$) & no \textbf{índice} \\
o passo lê & a célula: $\pi(t{+}1)=f(\pi(t))$ & o índice: $\pi(t{+}1)=g(t,\pi(t))$ \\
o campo $G$ lê & período e cauda (Teor.~\ref{thm:periodo}) & a dobra, e só ela \\
\bottomrule
\end{tabular}
\end{center}
O levantamento $\tilde\pi=(\pi,k)$ do Teor.~\ref{thm:levantamento} é a ponte
entre os dois --- e cobra o preço: junta espaço e história num só estado, e a
dobra desaparece, $\tilde G\equiv1$. As três coisas ---
\emph{dobrar}, \emph{ler só o estado} e \emph{não ter período} --- não podem
valer ao mesmo tempo.
\end{corolario}

\medskip\noindent\textbf{E não é uma opinião sobre a construção.} Que o passo
lê o índice verifica-se: procura-se na trajetória a mesma célula com sucessores
\emph{diferentes}. Na curva de Heighway ela existe --- a viragem
$\code{drag\_esq}(s)$ lê $s$, não lê onde a curva está --- e a testemunha exibe-se
com os quatro índices. No Julia, no relógio e na reta não existe.

\medskip\noindent\textbf{O relógio.} A rotação $J\colon(a,b)\mapsto(-b,a)$ é
função do estado e mais nada; é involutiva ao quadrado, $J^2=-\mathrm{id}$ e
$J^4=\mathrm{id}$. Pelo Teor.~\ref{thm:periodo} o campo lê-lhe o período sem
qualquer aparato: quatro células, cada uma com $G>1$, cauda zero. É o mesmo
relógio com que a obra conta as suas dobras, e a leitura ``uma colisão é uma
marca do contador; ler é seguir as marcas'' é a cláusula~(3) do
Teor.~\ref{thm:multiplicidade} dita noutras palavras.

% ═══════════════════════════════════════════════════════════════════════════
\section{A série do número $\pi$}\label{sec:serie}

\noindent Uma série é uma trajetória: a sucessão das suas somas parciais. Em
aritmética \emph{inteira} --- sem vírgula, com divisão inteira --- convergir tem
um significado exato e finito, e o campo $G$ lê-o.

\begin{definicao}[a soma parcial como realização]\label{def:serie}
Dada uma sucessão de termos inteiros $(a_k)$, ponha-se $S_0=a_0$ e
$S_{t+1}=S_t+a_{t+1}$. A realização é $\pi(t)=S_t$, com $X=\Z$.
\end{definicao}

\begin{teorema}[o campo lê o custo da série]\label{thm:serie}
Suponha-se que existe $T$ com $a_k=0$ para todo $k>T$ (os termos anulam-se em
tempo finito, que é o que a divisão inteira faz), e que
$S_0,\dots,S_T$ são distintos. Se $N\ge T+1$, então
\[
\bigl\lvert\{x:G(x)>1\}\bigr\rvert=1,\qquad
\bigl\lvert\{x:G(x)=1\}\bigr\rvert=T,
\]
isto é: o período é $1$ --- a trajetória entrou num ponto fixo --- e a cauda
conta exatamente as somas parciais distintas, que é o \textbf{custo} da série.
\end{teorema}

\begin{proof}
Para $t>T$ tem-se $S_{t+1}=S_t$, logo $\pi(t)=S_T$ para todo $t\ge T$: a
trajetória é constante a partir de $T$, o que é ser periódica de período $p=1$.
Como $S_0,\dots,S_T$ são distintos, o primeiro índice repetido é $\mu=T$, e o
Teor.~\ref{thm:periodo} aplica-se (a hipótese $N\ge\mu+2p-1=T+1$ é a do
enunciado). Donde $\lvert\{G>1\}\rvert=p=1$ e $\lvert\{G=1\}\rvert=\mu=T$.
\end{proof}

\medskip\noindent\textbf{O valor, e por que não é citado.} Com
$\arctan\frac1n=\sum_{k\ge0}\frac{(-1)^k}{(2k+1)n^{2k+1}}$ avaliado em escala
inteira $S=10^{8}$, três combinações independentes,
\[
16\arctan\tfrac15-4\arctan\tfrac1{239},\qquad
4\arctan\tfrac12+4\arctan\tfrac13,\qquad
8\arctan\tfrac12-4\arctan\tfrac17,
\]
têm de dar o mesmo número. Nenhum dígito é escrito à mão: se fossem copiados, a
igualdade não poderia falhar, e assim pode --- basta que qualquer uma das três
somas esteja errada. O número $\pi$ está \emph{dentro}: é produzido pela mesma
aritmética inteira que constrói o resto, e é sobre a trajetória das suas somas
parciais que o Teor.~\ref{thm:serie} corre.

\medskip\noindent\textbf{Onde ela cai nas duas dependências.} O termo $a_k$ é
função de $k$, não da soma corrente: a série, como trajetória, lê o
\emph{índice}. Ela dobra apenas no fim, ao entrar no ponto fixo --- e é aí, e só
aí, que passa a ser função do estado ($x\mapsto x$). É o Cor.~\ref{cor:duais} no
caso mais simples possível: a convergência \emph{é} a passagem de ler o tempo
para ler o espaço.

% ═══════════════════════════════════════════════════════════════════════════
\section{O dragão: a régua de bits e a dobra}\label{sec:dragao}

\noindent A curva de Heighway é a realização em $\Z^2$ que dobra mais
economicamente, e é onde $G>1$ se vê a olho. Ela admite duas construções sem uma
linha em comum, e é essa coincidência que a fixa.

\begin{definicao}[régua de bits]\label{def:bits}
Para $k\ge1$ escreva-se $k=2^{a}m$ com $m$ ímpar. Ponha-se
\[
\tau(k)=\begin{cases}+1 & m\equiv1\pmod 4\\[2pt] -1 & m\equiv3\pmod 4.\end{cases}
\]
A \textbf{curva de bits} $C_K$ é a sucessão de pontos $P_0=(0,0)$,
$P_{s+1}=P_s+u_s$, onde $u_0=(1,0)$ e $u_{s}=R_{\tau(s)}\,u_{s-1}$, com
$R_{+1}$ a rotação de $+90^{\circ}$ e $R_{-1}$ a de $-90^{\circ}$, para
$s=1,\dots,2^{K}-1$.
\end{definicao}

\begin{definicao}[a dobra]\label{def:dobrar}
Dada uma sucessão finita de pontos $D=(P_0,\dots,P_M)$ com extremo $F=P_M$, o seu
\textbf{dual} $D^{*}$ é
\[
D^{*}=\bigl(Q_0,\dots,Q_M\bigr),\qquad Q_j=F+R\,(F-P_{M-j}),
\]
com $R$ uma rotação fixa de $\pm90^{\circ}$. A \textbf{soma} $D+D^{*}$ é a
concatenação de $D$ com $D^{*}$ suprimindo o ponto repetido $Q_0=F$. Define-se
$D_0=\bigl((0,0),(1,0)\bigr)$ e $D_{k+1}=D_k+D_k^{*}$.
\end{definicao}

\noindent Em palavras: $D^{*}$ é a curva percorrida ao contrário e rodada de
$90^{\circ}$ em torno do ponto final. É a dobra do papel --- dobrar a tira ao
meio é colar-lhe a sua própria imagem invertida --- e é o mesmo passo
$T_{k+1}=T_k+T_k^{*}$ com que a obra constrói as suas torres, aqui numa curva.

\begin{lema}[antissimetria da régua]\label{lem:anti}
Para $1\le j<2^{k}$ vale $\tau(2^{k}+j)=-\tau(2^{k}-j)$.
\end{lema}

\begin{proof}
Escreva-se $j=2^{a}m$ com $m$ ímpar; como $j<2^{k}$, tem-se $a<k$. Então
\[
2^{k}-j=2^{a}\bigl(2^{k-a}-m\bigr),\qquad
2^{k}+j=2^{a}\bigl(2^{k-a}+m\bigr),
\]
e como $a<k$ o fator $2^{k-a}$ é par, logo $2^{k-a}\mp m$ são \emph{ímpares} ---
são portanto as partes ímpares de $2^{k}\mp j$. Ora
\[
\bigl(2^{k-a}-m\bigr)+\bigl(2^{k-a}+m\bigr)=2^{\,k-a+1}\equiv0\pmod 4,
\]
porque $k-a+1\ge2$. Dois ímpares cuja soma é $\equiv0\pmod4$ são um $\equiv1$ e
o outro $\equiv3$ módulo $4$. Logo $\tau$ toma valores opostos neles.
\end{proof}

\begin{teorema}[os dois caminhos dão a mesma curva]\label{thm:dragao}
Para todo $k\ge0$, a curva de bits $C_k$ da Def.~\ref{def:bits} e a curva
$D_k$ da Def.~\ref{def:dobrar} coincidem ponto a ponto, para uma das duas
escolhas do sinal de $R$ --- a mesma para todo $k$.
\end{teorema}

\begin{proof}
Indução em $k$. Para $k=0$ ambas são $\bigl((0,0),(1,0)\bigr)$.

Suponha-se $C_k=D_k$, com $M=2^{k}$ passos e extremo $F=P_M$. A curva $C_{k+1}$
tem $2^{k+1}$ passos e as suas primeiras $M$ direções são as de $C_k$ --- porque
$\tau(s)$ não depende de $k$ --- logo $C_{k+1}$ prolonga $C_k$ e o seu ponto de
índice $M$ é $F$. Resta identificar a segunda metade.

Para $1\le j<M$, a direção do passo $M+j$ de $C_{k+1}$ obtém-se da anterior pela
rotação $R_{\tau(M+j)}$, e pelo Lema~\ref{lem:anti} $\tau(M+j)=-\tau(M-j)$. Ou
seja: \emph{a sucessão de viragens da segunda metade é a da primeira metade lida
de trás para a frente e com os sinais trocados.}

Traduza-se isto em pontos. Percorrer $C_k$ ao contrário troca cada viragem pela
oposta --- é a definição de percorrer ao contrário, pois uma viragem à esquerda
vista do outro sentido é uma viragem à direita. Logo a sucessão de viragens de
$\operatorname{rev}(C_k)$ é exatamente $\bigl(-\tau(M-j)\bigr)_j$, que é a da
segunda metade de $C_{k+1}$. Duas curvas com a mesma sucessão de viragens e o
mesmo passo unitário diferem por uma isometria direta do plano; como ambas
começam em $F$, a isometria é a rotação em torno de $F$ que leva a primeira
direção de $\operatorname{rev}(C_k)$ na primeira direção da segunda metade de
$C_{k+1}$. Essa rotação é de $\pm90^{\circ}$, porque a direção do passo $M+1$ se
obtém da do passo $M$ por $R_{\tau(M+1)}$ e $\tau$ toma valores em
$\{\pm1\}$. Portanto a segunda metade de $C_{k+1}$ é
$F+R\,(F-P_{M-j})$, que é $D_k^{*}$, e $C_{k+1}=D_k+D_k^{*}=D_{k+1}$.

Que o sinal de $R$ é o \emph{mesmo} para todo $k$: ele é determinado por
$\tau(M+1)=\tau(2^{k}+1)$, e $2^{k}+1$ é ímpar com $2^{k}+1\equiv1\pmod4$ para
$k\ge2$, valor constante.
\end{proof}

\medskip\noindent\textbf{Por que isto é uma verificação e não uma repetição.} A
Def.~\ref{def:bits} decide cada viragem por uma conta sobre os bits do índice; a
Def.~\ref{def:dobrar} nunca olha para um índice --- copia, inverte e roda a curva
inteira. Não partilham uma linha. Quando duas construções assim coincidem em
$2^{12}$ pontos, o que se testou foi o objeto, não o contador.

\medskip\noindent\textbf{E o que não serve como verificação.} Comparar
$\abs{I}$ ou $\sum_x G(x)$ entre duas implementações \emph{não} distingue o
dragão de coisa nenhuma: pela cláusula (4) do Teorema~\ref{thm:multiplicidade},
$\sum G=\abs{I}$ vale para \emph{qualquer} realização com o mesmo número de
passos. Uma espiral quadrada passa nesse teste. O invariante tem de separar a
classe que se afirma estar a separar.

% ═══════════════════════════════════════════════════════════════════════════
\section{Realizações}\label{sec:realizacoes}

\noindent O Algoritmo~A corre em todas elas sem uma linha diferente.

\begin{center}\small
\begin{tabular}{@{}llp{0.44\textwidth}@{}}
\toprule
$n$ & realização & o que é \\
\midrule
$1$ & Cantor & o endereço ternário do terço médio; $G\equiv1$ --- o pó são
endereços distintos, e é por isso que tem medida nula sem se auto-intersetar \\
$2$ & Heighway & a curva da §\ref{sec:dragao}; $\max G=2$ \\
$3$ & dragão no espaço & a mesma regra de viragem com o \emph{plano} da viragem
a rodar de eixo; é uma realização em $\Z^3$, não ``o mesmo dragão com mais
espaço'' \\
$2$ & Julia sobre $\Z[i]$ & $z\mapsto z^{2}+c$ com tudo inteiro: determinista, e
o Teorema~\ref{thm:periodo} aplica-se \\
$2$ & o relógio $J$ & $(a,b)\mapsto(-b,a)$: função do estado, $J^4=\mathrm{id}$;
o campo lê-lhe o período $4$ com cauda $0$ \\
$1$ & as somas parciais & §\ref{sec:serie}: a série do número $\pi$ em inteiros;
período $1$, e a cauda é o custo \\
$4$ & o passo da máquina & $(\text{registo},\text{slot},\text{banco},\text{andar})$;
a aranha não sabe que é uma máquina, lê e escreve o mesmo $G$ \\
$k$ & a árvore do encaixe & §\ref{sec:ultra}: a célula \emph{é} a bola \\
\bottomrule
\end{tabular}
\end{center}

\medskip\noindent\textbf{Controlos.} Uma afirmação ``$G>1$ aqui'' só tem
conteúdo com uma realização onde $G\equiv1$ ao lado; senão passaria com um campo
que contasse mal. Servem a reta (injetiva por construção) e o Cantor. E note-se
que são \emph{dois}: o Cantor não dobra, apesar de ser a mais intrincada das
figuras da lista.

\medskip\noindent\textbf{$\max G$ é uma medida da curva.} Medido, o dragão de
$\Z^3$ da tabela tem $\max G=3$ e o do plano $2$: dobra \emph{mais}, e não menos.
A grandeza compara as duas curvas, que são construções distintas --- a de $\Z^3$
roda o plano da viragem --- e não os dois espaços. A dimensão, essa, já apareceu
onde tinha de aparecer: no Teor.~\ref{thm:dimensao}, $\lvert V(x)\rvert=2n$.

\medskip\noindent\textbf{E a memória é do algoritmo, não da tabela.} O
Teor.~\ref{thm:multiplicidade}(3) escreve num ponto por passo, e o suporte de $G$
tem no máximo $\abs{I}$ pontos: o custo é $\Theta(\abs{I})$ em qualquer $n$, como
a §\ref{sec:algoritmo} conta. Um arranjo denso sobre $X$ custaria
$\lvert X\rvert$ --- em $\Z^3$ com lado $2048$, $2048^3$ células --- e essa é uma
propriedade da tabela escolhida.

% ═══════════════════════════════════════════════════════════════════════════
\section{Medido}

\medido{Teor.~\ref{thm:multiplicidade} e Teor.~\ref{thm:dimensao}:
\code{tests/aranha\_n.c} §AN1--§AN2 (campo incremental $=$ recontagem e
$\sum G=\abs{I}$ em $\Z^1$--$\Z^4$), §AN7 ($\abs{V(x)}=2n$, vizinhos distintos);
Teor.~\ref{thm:dragao} §AN3 (bits $=$ dobra $D_k+D_k^{*}$, $k=1..12$, com os
dois sinais de $R$ corridos e exatamente um a bater);
dobra e controlos §AN4;
Teor.~\ref{thm:bolas} §AN6 (bolas particionam; todo ponto é centro; a euclidiana
falha ambas com o mesmo raio);
Teor.~\ref{thm:periodo} §AN8 ($\lvert\{G>1\}\rvert=p$ e $\lvert\{G=1\}\rvert=\mu$
sobre $z\mapsto z^2+c$ em $\Z[i]$, com o período obtido por busca independente e
a órbita que escapa como controlo);
Teor.~\ref{thm:leitura} e Cor.~\ref{cor:duais} §AN10 (a mesma célula do dragão
com sucessores diferentes, exibida; Julia, relógio e reta a lerem o estado; e o
período $4$ do relógio no campo);
Teor.~\ref{thm:serie} §AN11 (período $1$, cauda $=$ custo, e Machin, Euler e
Hermann a coincidirem sem um dígito escrito à mão);
Teor.~\ref{thm:levantamento} §AN5 e §AN9 ($\tilde G\equiv1$ pela própria aranha em
dimensão $n+1$; $\operatorname{pr}_1\circ\tilde\pi=\pi$ com resíduo $0$; fibra
vertical $\{1..G\}$; e a volta a devolver $G$ célula a célula);
o plano \code{tests/aranha\_g.c} §AG1--§AG11 e
\code{tests/aranha\_inversa.c} §AI1--§AI7.}

\colunas{Teor.~\ref{thm:multiplicidade} (só igualdade); Teor.~\ref{thm:levantamento}
(a fibra volta como uma coordenada); Teor.~\ref{thm:bolas} (partição);
Teor.~\ref{thm:periodo} ($\mu$ e $p$ no campo);
Teor.~\ref{thm:leitura} (dobrar sem período obriga a ler o índice)}
        {campo esparso, $O(1)$/passo, sem alocação dinâmica;
         $\Z^1$--$\Z^4$, árvore, $\Z[i]$}
        {\code{tests/aranha\_n.c} §AN1--§AN11;
         \code{tests/aranha\_g.c} §AG1--§AG11;
         \code{tests/aranha\_inversa.c} §AI1--§AI7}

\vfill
{\footnotesize\color{cinza}\noindent
Documento autocontido: todo enunciado usado está demonstrado acima.
Números: \code{bash tools/bateria.sh}.
Proprietário --- uso mediante contrato e pagamento (ver \texttt{LICENSE}).\par}

\end{document}
